% 底本：EGAIV.1.pdf，PDF第212--216页（分册印刷第213--217页）。
\Needspace{14\baselineskip}
\egasection{\S\ 23. 日本环}
\markright{\S\,23\quad 日本环}
\source{第四卷第一部分，原分册第213--221页；本节的23.1为第213--217页}
本节的结果将在第IV章（7.6）和（7.7）中得到补充。

\subsec{23.1. 日本环}
\egas{23.1.1}{定义}
称一个整环 $A$ 为\emph{日本环}，如果对其分式域 $K$ 的每个有限扩张 $K'$，
$A$ 在 $K'$ 中的整闭包 $A'$ 都是有限型 $A$-模（换言之，是有限 $A$-代数）。
称一个环 $A$ 为\emph{普遍日本环}，如果每个整的有限型 $A$-代数都是日本环。

若 $A$ 是日本环，$S\subset A$ 是不含 $0$ 的乘法集，则分式环 $S^{-1}A$ 也是日本环。
事实上，沿用上述记号，$S^{-1}A'$ 是 $S^{-1}A$ 在 $K'$ 中的整闭包，而且是有限型 $S^{-1}A$-模。
\cmt{P16-N01}{原文写“每个分式环”；按本书允许 $0\in S$ 的乘法集约定，须在这里排除这种情形，因为零环不是本定义中的整环。这里“有限”是模有限，不只是代数有限型；日本环定义本身没有加入诺特性。}

一个诺特整环，甚至一个离散赋值环，都不一定是日本环［30］。

\egas{23.1.2}{命题}
设 $A$ 是诺特整环，$K$ 是其分式域。若对 $K$ 的每个有限纯不可分扩张 $K'$，
$A$ 在 $K'$ 中的整闭包都是有限型 $A$-模，则 $A$ 是日本环。

\eproof
由于 $A$ 诺特，为验证 $A$ 是日本环，只须证明：对 $K$ 的每个有限拟Galois扩张 $L$，
$A$ 在 $L$ 中的整闭包 $B$ 是有限型 $A$-模。
令 $K'$ 为 $L$ 中包含的最大的 $K$ 的纯不可分扩张，则 $L/K'$ 为Galois扩张。
若 $A'$ 为 $A$ 在 $K'$ 中的整闭包，$B$ 也是 $A'$ 在 $L$ 中的整闭包。
根据假设，$A'$ 有限于 $A$，因而诺特且整闭。
已知在 $K'$ 的有限可分扩张中，$A'$ 的整闭包是有限型 $A'$-模
（Bourbaki，《交换代数》第V章，\S1，第6目，命题20的推论1）。故 $B$ 也有限于 $A$。
\finishproof

由（23.1.2）可知，当 $K$ 的特征为零时，诺特整环 $A$ 是日本环，当且仅当
它在自身分式域中的整闭包 $A'$ 是有限型 $A$-模。

\egas{23.1.3}{定理（Tate）}
设 $A$ 为诺特整环，$x\in A$ 且 $x\ne0$。假设：
\begin{enumerate}[label=(\roman*)]
\item $A$ 整闭；
\item $\mathfrak p=xA$ 为素理想，且 $A$ 对 $\mathfrak p$-预进制拓扑分离且完备；
\item $A/xA$ 为日本环。
\end{enumerate}
则 $A$ 为日本环。

我们将使用下述引理。

\egas{23.1.3.1}{引理}
设 $A$ 为环，$x\in A$ 不是 $A$ 中的零因子，且 $\mathfrak p=xA$ 为素理想。
则对每个整数 $n>0$，$x^nA_{\mathfrak p}$ 在 $A$ 中的逆像为 $x^nA$。

\eproof
设 $b\in A$，且在 $A_{\mathfrak p}$ 中有 $b/1=x^na/s$，其中 $a\in A$、$s\notin\mathfrak p$。
于是存在 $s'\notin\mathfrak p$，使 $s'sb=x^nas'$。
因此 $s'sb\in\mathfrak p$；由 $s's\notin\mathfrak p$ 得 $b\in\mathfrak p$，可写为 $b=xb'$，$b'\in A$。
因为 $x$ 不是零因子，可以消去 $x$，得到 $s'sb'=x^{n-1}as'$。
对 $n$ 归纳即得结论；逆向包含由定义给出。
\finishproof

\eproof
回到定理。由于 $A$ 整闭，可以只考虑其分式域 $K$ 的特征为 $p>0$ 的情形。
取 $K$ 的有限纯不可分扩张 $K'$，并选取 $q=p^e$ 使 $(K')^q\subset K$。
必要时扩大这一有限纯不可分扩张，还可假定存在 $y\in K'$ 使 $y^q=x$。
若扩大后整闭包有限，则原整闭包作为有限模的子模也有限，因为 $A$ 诺特。
令 $A'$ 为 $A$ 在 $K'$ 中的整闭包。由 $A$ 整闭，
\[
 A'=\{z\in K'\mid z^q\in A\}.
\]
令 $V=A_{\mathfrak p}$。其极大理想 $\mathfrak m=xV$ 为主理想，故这个诺特局部环是离散赋值环（17.1.4）。
与上面相同，$V$ 在 $K'$ 中的整闭包 $V'$ 为
\[
 V'=\{z\in K'\mid z^q\in V\}.
\]
已知 $V'$ 是离散赋值环，其极大理想及剩余域满足
\[
 \mathfrak m'=\{z\in K'\mid z^q\in\mathfrak m\},
 \qquad [V'/\mathfrak m':V/\mathfrak m]<\infty
\]
（Bourbaki，《交换代数》第VI章，\S8，第6目，命题6；第V章，\S2，第3目，引理4；
第VI章，\S8，第1目，引理2）。
\cmt{P16-T01}{依作者《勘误与增补（清单三）》A部，原分册第215页第4行的 $x'^q\in K'$ 改为 $x'\in K'$；同页第9行的 $x'^q\in A'$ 改为 $x'^q\in A$。下文按修正后的元素范围书写。}

对每个整数 $n>0$，我们证明
\begin{equation}
 (\mathfrak m')^n\cap A'=y^nA'.\tag{23.1.3.2}
\end{equation}
由 $y^q=x$ 得 $y\in\mathfrak m'$，故右端包含于左端。
反之，设 $x'\in(\mathfrak m')^n\cap A'$，写成 $x'=z'y^n$，其中 $z'\in K'$。
有 $x'^q\in A$。另一方面，$x'$ 是若干乘积 $t'_1\cdots t'_n$ 的和，其中每个 $t'_i\in\mathfrak m'$。
由于 $t_i'^q\in\mathfrak m$，而 $q$ 次幂在特征 $p$ 下保持加法，故
$x'^q\in\mathfrak m^n\cap A$。
引理（23.1.3.1）给出 $\mathfrak m^n\cap A=x^nA$，于是
$z'^qx^n\in x^nA$。因 $x\ne0$，可消去 $x^n$，得到 $z'^q\in A$，即 $z'\in A'$。
这证明了（23.1.3.2）。

其次，$A'$ 的 $xA'$-预进制拓扑是分离的。
该拓扑也就是 $yA'$-预进制拓扑；由（23.1.3.2），它是 $V'$ 的 $\mathfrak m'$-预进制拓扑所诱导的拓扑。
后者分离，因为 $V'$ 是离散赋值环。

再证明 $A'/xA'$ 是有限型 $A$-模。由于 $x=y^q$，且
$y^kA'/y^{k+1}A'\simeq A'/yA'$，只须证明 $A'/yA'$ 有限于 $A$。
在（23.1.3.2）中取 $n=1$，可知 $A'/yA'$ 嵌入域 $V'/\mathfrak m'$。
后者是 $V/\mathfrak m$ 的有限扩张，而 $V/\mathfrak m$ 为 $A/\mathfrak p$ 的分式域。
又因 $A'$ 整于 $A$，$A'/yA'$ 整于 $A/\mathfrak p$，所以包含在
$A/\mathfrak p$ 于 $V'/\mathfrak m'$ 中的整闭包内。
根据假设，$A/\mathfrak p$ 是诺特日本环，这个整闭包有限于 $A/\mathfrak p$；
因此其子模 $A'/yA'$ 有限于 $A/\mathfrak p$，从而有限于 $A$。

由于 $A'$ 的 $xA'$-预进制拓扑分离，$A'$ 嵌入其完备化 $\widehat{A'}$。
由 $A'/xA'$ 有限于 $A$，且 $A$ 对 $xA$-预进制拓扑分离完备，
（$0_{\mathrm I}$，7.2.9）说明 $\widehat{A'}$ 为有限型 $A$-模。
$A$ 诺特，故其子模 $A'$ 也有限型，定理得证。
\cmt{P16-T02}{原证明最后写“因此 $A$ 也如此”，应为 $A'$；否则没有得到所需整闭包的有限性。这里补明使用的是诺特环上有限模的子模仍有限生成。此项是译者校订，不列入上一注的作者正式勘误。}
\finishproof

\egas{23.1.4}{推论}
设 $A$ 是诺特且整闭的日本环，则每个形式幂级数环 $A[[T_1,\ldots,T_n]]$ 都是日本环。

\eproof
已知对诺特整闭环 $B$，$B[[T]]$ 仍诺特且整闭
（Bourbaki，《交换代数》第V章，\S1，第4目，命题14）。
因而对 $n$ 归纳，可归结为证明 $A[[T]]$ 是日本环。
取 $x=T$，它满足（23.1.3）的全部条件，即得结论。
\finishproof

\egas{23.1.5}{定理（Nagata）}
每个完备诺特局部整环 $A$ 都是日本环。

\eproof
由（19.8.8，ii），$A$ 含有一个完备正则局部子环 $B$，使 $A$ 为有限 $B$-代数。
只须证明 $B$ 是日本环：对 $\operatorname{Frac}(A)$ 的有限扩张，两环在该扩张中的整闭包相同，
且 $\operatorname{Frac}(A)$ 有限于 $\operatorname{Frac}(B)$。
因此可假设 $A$ 本身正则。

对 $n=\dim A$ 归纳。$n=0$ 时 $A$ 为域，结论成立。
设 $n>0$，令 $\mathfrak m$ 为 $A$ 的极大理想，取 $x\in\mathfrak m\setminus\mathfrak m^2$。
由（17.1.8），$A/xA$ 正则；由（17.1.7）及（16.3.4），其维数为 $n-1$。
此外，理想 $xA$ 闭（$0_{\mathrm I}$，7.3.5），故 $A/xA$ 完备。
归纳假设说明它是日本环。
$A$ 也对 $xA$-预进制拓扑分离且完备，所以（23.1.3）适用，得到 $A$ 为日本环。
\cmt{P16-G01}{原证明直接调用Tate定理。本处补明所需拓扑条件：$x^nA\subseteq\mathfrak m^n$，故一个 $x$-进制相容系给出 $\mathfrak m$-进制柯西列；各 $x^nA$ 在 $\mathfrak m$-进制拓扑下闭，其极限保持每阶同余。再由 $\bigcap\mathfrak m^n=0$ 得分离性。报告给出完整验证；这是补足应用细节，不是否定定理。}
\finishproof

\egas{23.1.6}{推论}
设 $A$ 为完备诺特局部整环，$K$ 为其分式域，$K'/K$ 为有限扩张。
则 $A$ 在 $K'$ 中的整闭包 $A'$ 是完备局部环。

\eproof
（23.1.5）已经说明 $A'$ 为有限 $A$-代数，因而是完备诺特半局部环，
于是为若干局部环的直积。由于 $A'$ 是整环，此直积只能有一个非零因子，所以 $A'$ 为局部环。
\finishproof

\egas{23.1.7}{命题}
设 $A$ 为诺特局部整环，$K$ 为其分式域，$K'/K$ 为有限扩张，$A'$ 为 $A$ 在 $K'$ 中的整闭包。
假设完备化 $\widehat A$ 约化，以 $R$ 表示 $\widehat A$ 的全分式环。

\noindent（i）若 $R\otimes_KK'$ 约化，则 $A'$ 为有限型 $A$-模。
特别地，当 $K'/K$ 可分时，此条件成立，因为 $R$ 是若干 $K$ 的扩域的直积
（Bourbaki，《代数》第VIII章，\S7，第3目，定理1的推论1）。

\noindent（ii）特别地，若 $R$ 为可分 $K$-代数，则 $A$ 为日本环。

\eproof
（i）简记 $A_1=\widehat A$，$A'_1=A'\otimes_AA_1$，$K'_1=K'\otimes_AA_1$。
由 $A_1$ 对 $A$ 平坦（$0_{\mathrm I}$，7.3.5），$A'_1$ 嵌入 $K'_1$，并且整于 $A_1$。
令 $K_1=K\otimes_AA_1$，则 $K'_1=K'\otimes_KK_1$。
由于 $A$ 为整环且 $A_1$ 平坦，$A$ 中每个非零元素在 $A_1$ 上都是正则的（$0_{\mathrm I}$，6.3.4）。
令 $S=A\setminus\{0\}$，于是 $K_1=S^{-1}A_1$ 可视为 $R$ 的子环。
又因每个 $K$-模都平坦，$K'_1$ 可视为 $K'\otimes_KR$ 的子环；后者按假设约化，且有限于 $R$。

设 $\mathfrak q_i$（$1\le i\le n$）为 $A_1$ 的极小素理想，令
$B_i=A_1/\mathfrak q_i$，$L_i=\operatorname{Frac}(B_i)$。
则 $R=\prod_iL_i$，而 $K'\otimes_KR=\prod_i(K'\otimes_KL_i)$。
每个 $K'\otimes_KL_i$ 都是约化的有限 $L_i$-代数，因此为 $L_i$ 的有限扩域的有限直积。
$B_i$ 是完备局部整环；由Nagata定理（23.1.5），$B_i$ 在这个直积中的整闭包
是有限型 $B_i$-模，从而也是有限型 $A_1$-模。
每个因子中的元素对 $A_1$ 整，当且仅当对其像 $B_i$ 整，所以
$A_1$ 在 $K'\otimes_KR$ 中的整闭包是有限型 $A_1$-模。
$A'_1$ 包含在这个整闭包中，$A_1$ 又诺特，故 $A'_1$ 有限型。
最后，$A_1$ 对 $A$ 忠实平坦（$0_{\mathrm I}$，7.3.5），由有限型模的忠实平坦下降
（Bourbaki，《交换代数》第I章，\S3，第6目，命题11），$A'$ 有限于 $A$。

（ii）假设说明各 $L_i/K$ 可分，因此 $K'\otimes_KL_i$ 约化
（Bourbaki，《代数》第VIII章，\S7，第3目，定理1）。
对每个有限扩张 $K'/K$ 应用（i），即得结论。
\finishproof
